% equilibrium.tex — Paper 6, Section IV: Entanglement Equilibrium
% ASSERT_CONVENTION: natural_units=natural, metric_signature=mostly_plus, entropy_base=nats, modular_hamiltonian=K_minus_ln_rho, coupling_convention=H_sum_hxy, state_normalization=Tr_rho_1
\section{Entanglement Equilibrium and the Thermal Time Hypothesis}
\label{sec:equilibrium}

The Jacobson entanglement equilibrium argument~\cite{Jacobson2016}
requires a third input beyond area-law entanglement and the
entanglement first law: the maximal vacuum entanglement hypothesis
(MVEH), the condition that the vacuum maximizes entanglement entropy at
fixed stress-energy expectation value.  In the original Jacobson
framework, MVEH is an additional assumption.  We argue that in the
present context---where there is no pre-existing geometry---it is
\emph{definitional}.

\subsection{Modular flow and the thermal time hypothesis}
\label{sec:thermal-time}

The Tomita--Takesaki theorem guarantees that every faithful normal
state $\omega$ on a von~Neumann algebra $\mathcal{M}$ defines a
canonical one-parameter automorphism group
$\sigma_t^\omega : \mathcal{M} \to \mathcal{M}$, the \emph{modular
automorphism group} (or modular flow)~\cite{BratteliRobinson1997}.
The flow is determined entirely by the pair $(\mathcal{M}, \omega)$:
no external Hamiltonian or time parameter is needed.

Connes and Rovelli~\cite{ConnesRovelli1994} proposed the \emph{thermal
time hypothesis}: in a generally covariant quantum theory---one with
no preferred notion of time---the physical time evolution \emph{is}
the modular flow.  Time is not a background parameter; it is selected
by the state.  This proposal was motivated by the observation that the
modular flow of a thermal state $\omega_\beta$ at inverse temperature
$\beta$ is precisely the physical time evolution rescaled by $\beta$:
the thermal state ``knows'' its own dynamics.

For the self-modeling lattice, there is no pre-existing spacetime
geometry.  There is no background metric, no pre-assigned light cone,
no external time parameter.  The lattice carries an algebra
$\mathfrak{A}$ (the quasi-local algebra of
Sec.~\ref{sec:lattice-def}) and a state.  The modular flow of this
state is the \emph{only} candidate for a local time evolution in the
absence of external structure.

Near the boundary $\partial A$ of a subregion $A$, the modular flow
plays the role of a local boost generator.  In quantum field theory on
a fixed spacetime, this is the content of the Bisognano--Wichmann
theorem~\cite{BisognanoWichmann1976}: the modular Hamiltonian of the
Minkowski vacuum restricted to a Rindler wedge is the Lorentz boost
generator, and the modular flow is the boost.  Here, we \emph{define}
the emergent geometric structure from the modular flow, rather than
deriving the modular flow from a pre-existing geometry.

\subsection{MVEH as definitional}
\label{sec:mveh}

The maximal vacuum entanglement hypothesis states that the vacuum
maximizes $\ent(B)$ for any small geodesic ball $B$ at fixed
stress-energy expectation value $\langle\Tab\rangle$.  When a
background spacetime exists, this is a substantive physical
assumption: one must check that the vacuum of a given quantum field
theory actually satisfies this extremal condition.

In the self-modeling framework, however, there is no pre-existing
geometry against which to test the hypothesis.  The logical structure
is different:

\begin{quote}
In the absence of pre-existing geometry, the modular flow of the
self-modeling lattice state defines the emergent geometric flow
(thermal time hypothesis,
Connes--Rovelli~\cite{ConnesRovelli1994}).  The state whose
modular flow defines a smooth emergent geometry necessarily
satisfies MVEH---the ``vacuum'' \emph{is} the geometry-defining
state by definition.
\end{quote}

The vacuum is not a state that happens to maximize entanglement; it is
the state \emph{selected by the requirement} that its modular flow
define a smooth geometric flow.  Entanglement equilibrium ($\delta\ent = 0$
under perturbations that preserve $\langle\Tab\rangle$) is the
\emph{selection criterion} for which state defines the geometry, not
an additional assumption about a pre-selected state.

This parallels the dissolution of the involution in
Paper~5~\cite{Paper5}.  There, the involution $J$ on $M_n(\mathbb{C})^{sa}$
was not assumed and then matched to an external standard---$J$ was the
structure that self-modeling produced, and there was no external
standard to match against.  Here, the geometry is not assumed and then
tested for MVEH---the geometry is defined by the state's modular flow,
and MVEH is the defining condition.

MVEH does not appear as a numbered assumption in this paper.

\subsection{The Sorce caveat}
\label{sec:sorce}

Sorce~\cite{Sorce2024} proved that geometric modular flow---modular
flow that acts as a local geometric symmetry (a conformal Killing
flow)---requires conformal symmetry of the underlying quantum field
theory.  Non-isometric conformal Killing flows arise only in conformal
field theories.  This places a constraint on the thermal time
identification: not every modular flow corresponds to a smooth
geometric flow.

For the self-modeling lattice, this caveat is addressed by the
$SU(n)$~symmetry of the Hamiltonian~\eqref{eq:hamiltonian}.  The
diagonal $U(n)$~covariance, combined with the nearest-neighbor
structure, places the infrared description of the lattice in the class
of $SU(n)$~symmetric spin models.  For $n = 2$ with $J > 0$
(antiferromagnetic), the infrared theory is the $SU(2)_1$
Wess--Zumino--Witten conformal field theory with central charge
$c = 1$~\cite{Affleck1986, AffleckHaldane1987}.  For general~$n$ with
$J > 0$, the $SU(n)$ Heisenberg model flows to the $SU(n)_1$ WZW
CFT~\cite{Affleck1986b}.  In these conformal theories, the modular
Hamiltonian of the vacuum restricted to an interval takes the
Casini--Huerta--Myers form~\cite{CasiniHuertaMyers2011}, and the
Sorce condition is satisfied.

We state honestly that Sorce's result provides \emph{necessary}
conditions for geometric modular flow; \emph{sufficient} conditions
remain an open question.  The resolution via $SU(n)$ conformal
symmetry in the infrared is clean for the one-dimensional case
($d = 1$), where the WZW CFT description is well established.

For $d \geq 2$, this resolution faces a genuine open problem.  The
$SU(n)$ Heisenberg antiferromagnet on the square lattice ($d = 2$) is
expected to exhibit N\'eel order---spontaneous breaking of $SU(n)$
symmetry---rather than conformal behavior.  The infrared theory is
then a nonlinear sigma model with a mass gap (spin-wave gap), not a
CFT.  The conformal modular Hamiltonian
(Sec.~\ref{sec:conformal-approx}) would apply only in the window
$a \ll \ell \ll 1/m$ (where $m$ is the N\'eel gap scale), if such a
window exists.  This is the principal limitation of the present
approach in dimensions $d \geq 2$: the paper's results are strongest
in $d = 1$, where the WZW conformal structure is exact.

\subsection{Wilsonian continuum limit}
\label{sec:continuum}

At length scales $L \gg a$ (lattice spacing), we assume that the
lattice details become irrelevant by a Wilsonian universality
argument: the long-wavelength physics is governed by the
symmetry-allowed effective field theory, and the smooth Riemannian
manifold $(M, \gab)$ emerges at long wavelengths.

The Wilsonian continuum limit is standard \emph{methodology}, employed
in lattice QCD, causal dynamical triangulations, and Regge calculus.
However, its \emph{applicability} to a given lattice model is not
automatic.  Lattice QCD benefits from asymptotic freedom, which
guarantees a UV fixed point and controlled continuum limit.  The
self-modeling Hamiltonian has no such guarantee in $d \geq 2$: whether
the Heisenberg model on a $d$-dimensional lattice admits a continuum
limit with the correct universality class is an open question,
particularly for $d \geq 3$ where the model may be trivial or
confining.  A rigorous constructive proof is a hard open problem shared
by all lattice quantum gravity programs.  We frame the continuum limit
as a physical argument, not a rigorous construction.

The Wilsonian picture suggests a natural perspective from the
self-modeler:  the compressed model $M$ is finite-dimensional.  Below
the model's resolution (the lattice scale~$a$), there is nothing to
see.  Smoothness is not a property of the lattice; it is a property of
the observer's coarse-grained description at scales $L \gg a$.  The
continuum limit is the self-modeler's view.

In the continuum limit, the lattice parameters acquire physical
identifications:
\begin{itemize}
\item The lattice spacing $a$ maps to the Planck scale:
  $a \lesssim 2\ell_P \sqrt{\log n}$ (via the channel capacity
  bound~\eqref{eq:channel-bound} and
  $\Gnewton = 1/(4\eta)$).
\item The Lieb--Robinson velocity $\vlr$ maps to the speed of
  light~$c$.
\item The lattice entropy density
  $\eta_{\mathrm{lattice}} \leq \log n$ per boundary bond maps to
  the continuum entropy density
  $\eta = \eta_{\mathrm{lattice}}/a^{d-1}$.
\end{itemize}
